\scanpage{99}
Cas des triangles \uline{verticaux} $\{\alpha,s,s'\}$ $(\alpha\in t_0,s\in\widetilde E_\alpha,s'=\sigma_\alpha(s))$ On trouve

\uline{Prop 5} \quad Soient $\alpha\in\unsure{t|t_0}$, $s,s'$ deux arêtes liées de $E_\alpha$, et $\deleted{\beta\in t_0,\ \beta\ne\unsure{\beta|\alpha},}$ $t\in\widetilde E_\beta$, avec $\beta\in t_0$, $\beta\ne\alpha$, alors $t$ est lié à \deleted{au plus} $s$ ou à $s'$ (\uline{NB} par prop. 2, ou \unsure{rien que sous|si on veut par|encore suivant} la condition $(\alpha)$, \deleted{il est lié à un} il n'est pas lié à $s$ \uline{et} à $s'$, donc il est lié à un et un seul des deux pts $s,s'$). En termes des données du Scholie : Soit $\widetilde\Gamma_{\alpha,\beta}\subset\widetilde E_\alpha\times\widetilde E_\beta$ l'image de $\widetilde\Gamma$ dans $\widetilde E_\alpha\times\widetilde E_\beta$, soient $s,s'\in\widetilde E_\alpha$, $t\in\widetilde E_\beta$, avec $s'=\sigma_\alpha s$ i.e. $s=\sigma_\alpha s'$. Alors on a $\{s,t\}\in\unsure{\widetilde E_{\alpha\beta}|\widetilde\Gamma_{\alpha\beta}}$ ou $\{s',t\}\in\widetilde\Gamma_{\alpha\beta}$

\noindent\vrule\hspace{.5em}\begin{minipage}{.89\linewidth}\uline{Cor} \quad \deleted{\unsure{Alors}} Sous les conditions précédentes, posant $t'=\sigma_\beta t$, si \inserted{$s,t$ liés i.e.} $\deleted{(s,t)\in\unsure{\Gamma|\widetilde\Gamma}}\ \{s,t\}\in A$ i.e. $(s,t)\in\widetilde\Gamma_{\alpha\beta}$, alors \inserted{$s',t'$ liés i.e.} $\{s',t'\}\in A$ i.e. $(s',t')\in\widetilde\Gamma_{\alpha\beta}$.

\end{minipage}

En effet, $t'$ étant lié à $s$ ou $s'$, et ne pouvant être lié à $s$ (car $s$ serait lié à $t$ et $t'$ en opposition à prop 2) il est lié à $s'$.

\noindent\vrule\hspace{.5em}\begin{minipage}{.89\linewidth}\uline{Cor} \quad Soit $\deleted{\unsure{\sim}}\sigma_{\widetilde p}=\prod_{\alpha\in t_0}\deleted{\unsure{\alpha}}\overline\sigma_\alpha$ l'autom \inserted{\unsure{inv.}} d'ordre 2 (ss pt fixe) de $\prod\widetilde E_\alpha$. Alors $\widetilde\Gamma$ est stable par $\unsure{\ll|\sim}\sigma_{\widetilde p}\unsure{\gg|{>}}$\unsure{?|\textquoteright|;}\end{minipage}

On note $\sigma_{\widetilde\Gamma}$ l'autom d'ordre 2 de $\widetilde\Gamma$ induit par $\sigma_{\widetilde p}$.
% archive sheet number: 98
